% !TeX program = xelatex
\documentclass[11pt,a4paper]{article}

\usepackage[margin=22mm]{geometry}
\usepackage{amsmath,amssymb}
\usepackage{booktabs,longtable,array}
\usepackage{fontspec}
\usepackage{xeCJK}
\usepackage{iftex}
\usepackage[hidelinks]{hyperref}
\usepackage{xurl}

\ifPDFTeX
  \errmessage{This report contains CJK text and must be compiled with XeLaTeX}
\fi

\setmainfont{DejaVu Serif}
\setsansfont{DejaVu Sans}
\setmonofont{DejaVu Sans Mono}
\IfFontExistsTF{Noto Sans CJK SC}{
  \setCJKmainfont{Noto Sans CJK SC}
  \setCJKsansfont{Noto Sans CJK SC}
}{
  \IfFontExistsTF{Droid Sans Fallback}{
    \setCJKmainfont{Droid Sans Fallback}
    \setCJKsansfont{Droid Sans Fallback}
    \setCJKmonofont{Droid Sans Fallback}
  }{
    \setCJKmainfont{AR PL UMing CN}
    \setCJKsansfont{AR PL UMing CN}
    \setCJKmonofont{AR PL UMing CN}
  }
}

\title{Beam-Spline Rotational Recovery\\
\large SO(3) Finite-Rotation Hybrid、解析伴随与验证报告}
\author{}
\date{28 July 2026}

\newcommand{\RR}{\mathbb{R}}
\newcommand{\T}{\mathsf{T}}
\newcommand{\curl}{\nabla\times}
\newcommand{\relLtwo}{\mathrm{rel\_l2}}

\begin{document}
\maketitle

\begin{abstract}
本文审查并验证 Kratos MappingApplication 中
BeamSplineMapper 与
BeamSplineMapperWithRecoveryOfRotations。
核心结论是：Ahrem 等人的 generalized interpolation 是
small-rotation theory；当前 finite mode 是在其 rotation recovery
之上加入 $SO(3)$ rigid-section kinematics 的 hybrid extension，
并非 Ahrem 原文中的 finite theory。当前实现的 Rodrigues forward、
SO(3) analytical tangent 和 $J^\T f$ load transfer 已得到
analytical verification。Production 修改严格限制在两个 mapper 的
四个 \texttt{.h/.cpp} 文件；Kratos core、structural element、
solver 和 co-rotation mapper 均未修改。
\end{abstract}

\section{问题定义与四类正确性}

研究问题可表述为：
\begin{quote}
能否将 Ahrem rotational recovery 与 $SO(3)$ rigid-section
kinematics、analytical tangent-transpose load transfer 结合，
得到 finite-rotation、work-consistent 的 beam-to-surface mapper，
并明确它与 structural rotation parameterization 的兼容边界？
\end{quote}

必须区分以下四类结论：
\begin{enumerate}
  \item \textbf{Theoretical correctness}：运动学和离散 virtual work
        是否自洽；
  \item \textbf{Implementation correctness}：代码是否实现该运动学、
        tangent、符号及 state lifecycle；
  \item \textbf{Analytical benchmark correctness}：manufactured fields
        和离散测试是否达到规定误差；
  \item \textbf{FSI stability}：耦合系统能否在给定 nonlinear solver
        和 coupling parameters 下推进。FSI failure 不能自动推翻前三项。
\end{enumerate}

\section{$SO(3)$ 与 finite rigid-section kinematics}

\subsection{Rotation group 与 exponential coordinates}

三维 proper rotations 构成 special orthogonal group
\[
 SO(3)=\{R\in\RR^{3\times3}\mid R^\T R=I,\ \det R=1\}.
\]
$R^\T R=I$ 保持长度和夹角，$\det R=1$ 排除 reflection 并保持
handedness。Rotation vector 写成
\[
 \boldsymbol\theta=\phi\boldsymbol n,\qquad
 \phi=\|\boldsymbol\theta\|_2,\qquad \|\boldsymbol n\|_2=1.
\]
其 skew matrix $K=[\boldsymbol\theta]_\times$ 满足
$K v=\boldsymbol\theta\times v$。Exponential map 与 Rodrigues formula 为
\[
 R(\boldsymbol\theta)
 =\exp(K)
 =I+\frac{\sin\phi}{\phi}K
   +\frac{1-\cos\phi}{\phi^2}K^2.
\]
当 $\phi\rightarrow0$ 时，
\[
 \frac{\sin\phi}{\phi}=1-\frac{\phi^2}{6}+O(\phi^4),\qquad
 \frac{1-\cos\phi}{\phi^2}=\frac12-\frac{\phi^2}{24}+O(\phi^4),
\]
因此 $R=I+K+\frac12K^2+\cdots$；只保留 $I+K$ 才是
small-rotation approximation。

\subsection{Reference offset}

对 surface point $X_s^0$，先在 reference beam chain 上投影得到
$X_c^0$，并定义
\[
 r_0=X_s^0-X_c^0.
\]
Finite rigid-section displacement 是
\[
 u_s=u_c+\left(R(\theta_c)-I\right)r_0.
\]
$r_0$ 必须来自 reference configuration。若从 current coordinates
反推 offset，forward operator 会把未知 deformation 再次写入其几何，
造成 double counting、path dependence，并使 analytical tangent
不再对应实际 forward map。Mapper 也不得通过重置 current coordinates
来伪造 reference projection。

\section{Ahrem small theory 与当前 finite hybrid}

Ahrem rotational recovery 构造三维 vector field $s:\RR^3\to\RR^3$：
\[
 s(X_i)=u_i,\qquad
 (\curl s)(X_i)=2\theta_i.
\]
Factor 2 来自 infinitesimal kinematics
\[
 \theta=\frac12\curl u.
\]
其 small-mode surface map 是
\[
 u_s=s(X_s^0).
\]
这里 rotations 作为 derivative functionals 进入 generalized
RBF/Hermite interpolation；它并不要求某一截面上所有 points
严格保持刚性。

当前 finite hybrid 改为在 center projection 上读取 field：
\[
 u_c=s(X_c^0),\qquad
 \theta_c=\frac12(\curl s)(X_c^0),
\]
再用
\[
 u_s=u_c+\left[\exp([\theta_c]_\times)-I\right]r_0.
\]
因此二者区别不是 ``small formula 换成 Rodrigues'' 这么简单：
\begin{center}
\begin{tabular}{p{0.27\linewidth}p{0.31\linewidth}p{0.31\linewidth}}
\toprule
 & Ahrem small recovery & SO(3) rigid-section hybrid\\
\midrule
Evaluation point & 直接计算 $s(X_s^0)$ &
在 $X_c^0$ 计算 $s,\curl s$\\
Rotation meaning & infinitesimal curl &
rotation vector passed to $\exp:\mathfrak{so}(3)\to SO(3)$\\
Section constraint & 不强制完全刚性 &
reference cross-section 严格 rigid\\
Large rigid rotation & 一般不 exact &
rotation-vector contract 成立时 exact\\
Literature status & Ahrem 原始理论 &
当前工作的 finite extension\\
\bottomrule
\end{tabular}
\end{center}

Small/finite mode 不能在 Newton 或 coupling iteration 中按角度自动切换。
两者是不同 forward operators；切换会造成 residual 和 tangent
不连续，并使 load history/path 依赖于阈值。

\section{Analytical tangent 与 work consistency}

记 interpolation coefficients 为 $c$，离散系统为
\[
 Ac=Bq.
\]
$A$ 仅由 reference support、kernel 和 polynomial basis 决定。
$Bq$ 的 displacement rows 使用 $u_i$，curl rows 使用 $2\theta_i$。
对 destination point，finite forward 可写成
\[
 F(q)=E_c c+\bigl(R(\tfrac12C_c c)-I\bigr)r_0 .
\]
Right Jacobian 为
\[
 J_r(\theta)
 =I-\frac{1-\cos\phi}{\phi^2}[\theta]_\times
  +\frac{\phi-\sin\phi}{\phi^3}[\theta]_\times^2.
\]
Offset 对 rotation-vector variation 的导数是
\[
 \delta(Rr_0)
 =-R[r_0]_\times J_r(\theta)\,\delta\theta.
\]
因此 point tangent operator 为
\[
 T_c
 =E_c+\frac12\left[-R[r_0]_\times J_r(\theta_c)\right]C_c,
\qquad
 J=T_cA^{-1}B.
\]

对 surface load $f_s$，一致的 inverse map 必须满足
\[
 Q_b=J(q)^\T f_s,\qquad
 \delta q^\T Q_b=\delta u_s^\T f_s.
\]
实现不显式形成 $A^{-1}$ 或 $J$，而是先聚合
$T_c^\T f_s$，再求解
\[
 A^\T z=\sum_s T_{c,s}^\T f_s,\qquad Q_b=B^\T z.
\]
$B^\T$ 把 curl multipliers 映回 nodal moments，因此 rotations 的
factor 2 必须在此保留。Global sign 只由最后的
\texttt{SWAP\_SIGN} accumulation 控制。Finite difference 只用于
verification；若在 production 中逐 DOF perturb，不仅代价过高，
还会引入 step-size noise 和 coupling nondeterminism。

\section{三个 mapper 的模型边界}

\begin{longtable}{p{0.20\linewidth}p{0.24\linewidth}p{0.25\linewidth}p{0.22\linewidth}}
\toprule
 & Plain BeamSpline & Recovery mapper & BeamMapper co-rotation\\
\midrule
\endhead
Centerline model &
1D global chain spline；rotation 作为 slope/torsion data &
3D generalized value/curl interpolation &
moving/local frame 中的 beam interpolation\\
Surface offset &
由 interpolated section rotation 旋转 &
small: $s(X_s)$；finite: $SO(3)$ rigid section &
分解并组合 rigid-body operators\\
Large rotation &
含 finite matrix，但不是一般 objective/geometrically exact &
finite rigid rotation exact；受 rotation-vector contract 限制 &
通过 co-rotating frame 处理 large motion\\
Inverse &
当前 nonlinear forward tangent 的 transpose &
small linear adjoint；finite analytical $J^\T$ &
现有实现使用 section force/moment distribution；未在本工作中修改\\
\bottomrule
\end{longtable}

\section{CrBeamElement3D2N compatibility}

\texttt{CrBeamElement3D2N} 的 nodal \texttt{ROTATION} 是 additive
solution DOF，但 element 内部 orientation 通过 iteration increments
更新 quaternion。对同一固定轴或 commuting increments，
\[
 \exp(\theta_1)\exp(\theta_2)=\exp(\theta_1+\theta_2),
\]
所以 nodal total vector 与 accumulated orientation 可兼容。
一般情况下
\[
 \exp([\Delta\theta_2]_\times)\exp([\Delta\theta_1]_\times)
 \ne
 \exp([\Delta\theta_1+\Delta\theta_2]_\times).
\]
实测 single-axis small increments 的 orientation
$\relLtwo=6.8041\times10^{-14}$，通过 $10^{-10}$ gate；
non-commuting sequence 的差异为
$\relLtwo=7.4178\times10^{-3}$。

因此本实现保持 mapper 的明确 contract：
\texttt{ROTATION} 被解释为 global total exponential coordinates。
Mok planar/single-axis motion 属于兼容范围。对 general 3D
non-commuting CrBeam motion，只报告 limitation，不宣称 exact。
直接读取 element internal axes 并不能解决 load mapping，因为当前
四文件范围内无法取得该 orientation 对 nodal DOFs 的 consistent
derivative；这样做会破坏 $J^\T f$。

\section{Implementation changes}

Production 仅修改以下四个文件：
\begin{itemize}
  \item \texttt{beam\_spline\_mapper.h/.cpp};
  \item
  \texttt{beam\_spline\_mapper\_with\_recovery\_of\_rotations.h/.cpp}.
\end{itemize}

完成的 correctness hardening 为：
\begin{enumerate}
  \item plain mapper 和 recovery-finite mapper 的
        \texttt{InverseMap} 通常要求当前 interface 上已有成功
        \texttt{Map}；唯一例外是 standard
        \texttt{DISPLACEMENT/ROTATION} 均为零的无歧义初始态，
        用于兼容 CoSimulation 首次 load transfer；
  \item interface rebuild/update 会清除 forward state；
  \item finite cache 只在所有 destination nodes 成功完成 forward 后发布；
  \item \texttt{polynomial\_level=0} 才允许 auto selection；
        显式 level 1--4 若 node count/rank/unisolvency 不满足立即报错；
  \item legacy \texttt{polynomial\_basis} 保持兼容并发出
        deprecation warning；
  \item 保持原 reference offset、Rodrigues Taylor branch、right
        Jacobian、factor 2、transpose solve 和 \texttt{SWAP\_SIGN}。
\end{enumerate}

\section{Analytical verification results}

统一 vector/matrix error 为
\[
 \relLtwo(a,b)
 =\frac{\|a-b\|_2}{\max(\|b\|_2,\epsilon)}.
\]
Virtual work 另外报告 normalized scalar residual。

\subsection{Forward 与 polynomial span}

\begin{center}
\begin{tabular}{lcc}
\toprule
Test & Maximum $\relLtwo$ & Gate\\
\midrule
Finite rigid rotations $1^\circ,30^\circ,90^\circ,170^\circ$
  & $3.98\times10^{-15}$ & $10^{-10}$\\
Default small versus explicit small & $0$ & $10^{-12}$\\
Cubic-kernel rigid exact span, levels 0,2,3,4
  & $3.36\times10^{-16}$ & $10^{-10}$\\
Coordinate mutation & $0$ & exact zero\\
\bottomrule
\end{tabular}
\end{center}

Full-3D level 1 在 collinear support 上 rank deficient，按 strict
contract 正确拒绝。Gaussian kernel 的 level-2 rigid test 得到
$1.0384\times10^{-10}$，略高于 gate；同一 exact-span test 使用
cubic kernel 为 $1.7000\times10^{-16}$，说明这是 kernel/system
conditioning 而非 span 缺失。Out-of-span quadratic field 的误差约
$1.36\times10^{-1}$，只作为 approximation evidence，不标记为
exact-span failure。

\subsection{Tangent、inverse 与 signs}

\begin{center}
\begin{tabular}{lcc}
\toprule
Test & Worst result & Gate\\
\midrule
Finite analytical directional tangent vs central FD
  & $\relLtwo=1.5544\times10^{-9}$ & $10^{-7}$\\
Recovery-small virtual work & $2.2384\times10^{-9}$ & $10^{-8}$\\
Recovery-finite virtual work & $2.5748\times10^{-9}$ & $10^{-8}$\\
Plain spline virtual work & $2.3979\times10^{-9}$ & $10^{-8}$\\
\texttt{SWAP\_SIGN} & $\relLtwo=0$ & $10^{-12}$\\
\bottomrule
\end{tabular}
\end{center}

Surface Cartesian basis loads 被逐一 inverse mapped，用于重构
$Jp$；因此该 tangent test 同时覆盖 force distribution、
moment sign、rotation factor 2 和 transpose orientation。

\subsection{Bending、torsion 与 refinement}

使用 prescribed finite rigid-section reference，beam node counts 为
$6,9,17,33$：
\begin{center}
\begin{tabular}{lcc}
\toprule
Case & Observed $\relLtwo$ range & Status\\
\midrule
Pure bending & $6.16\times10^{-15}$--$1.32\times10^{-14}$ & PASS\\
Pure torsion & $2.43\times10^{-16}$--$4.29\times10^{-16}$ & PASS\\
Combined bending/torsion & $4.37\times10^{-15}$--$8.76\times10^{-15}$ & PASS\\
\bottomrule
\end{tabular}
\end{center}

\section{FSI gate 与 VTK evidence policy}

FSI 不参与上述 analytical acceptance。只有代码、tangent、
work consistency、CrBeam boundary 和本文档完成后才运行 Mok gate。
固定参数为
\[
 scale=50,\quad\Delta t=0.05,\quad\alpha=0.03,\quad
 N_{\rm coupling}=60,\quad r_k=0.50,\quad
 \lambda=10^{-8},\quad polynomial\_level=4.
\]
Cases 为 nearest-neighbor baseline、plain、recovery-small 和
recovery-finite；time gates 为
$0.5\rightarrow4.5\rightarrow10\rightarrow25$ s。

每次运行必须使用 versioned path，并生成 VTK manifest。Validation
检查 readable files、monotonic step/time、node count、NaN/Inf 和
required arrays；CSM 至少包括
\texttt{DISPLACEMENT}, \texttt{ROTATION}, \texttt{POINT\_LOAD},
\texttt{POINT\_MOMENT}, \texttt{REACTION},
\texttt{REACTION\_MOMENT}。监测节点的 VTK displacement 必须与
runtime history 用 $\relLtwo$ 交叉验证。异常终止时，只信任完整写出
且不晚于 accepted solution time 的 VTK。

\subsection{Gate results}

所有四种 case 均通过 0.5 s gate，且 VTK integrity、required arrays
和 input restoration 均通过。4.5 s gate 的结果为：
\begin{center}
\begin{tabular}{lcccc}
\toprule
Case & Status & Failure time & Last valid time & VTK\\
\midrule
Nearest neighbor & PASS & --- & 4.50 & valid\\
Plain spline & structural nonconvergence & 3.10 & 3.05 & valid\\
Recovery-small & structural nonconvergence & 4.25 & 4.20 & valid\\
Recovery-finite & structural nonconvergence & 3.10 & 3.05 & valid\\
\bottomrule
\end{tabular}
\end{center}
按 gate 规则，三个 spline cases 不进入 10 s。Nearest-neighbor 已有
充分 baseline evidence，因此不继续重复长时运行。这里的 ``VTK
valid'' 仅指截至 last accepted time 的完整文件可信，不能把 failed
step 当作已 flush。

\subsection{Nested convergence conditions}

该计算有两个不能混淆的 convergence gate。Structural solver 因为
存在 rotation DOFs，实际使用
\texttt{ResidualDisplacementAndOtherDoFCriteria}：
\[
 \left(r_u\le10^{-6}\ \lor\ a_u<10^{-6}\right)
 \land
 \left(r_\theta\le10^{-6}\ \lor\ a_\theta<10^{-6}\right),
\]
并限制为 150 Newton iterations。Coupling solver 使用 fluid-interface
displacement residual $\rho=d^{(k)}-d^{(k-1)}$：
\[
 a_c=\frac{\|\rho\|_2}{\sqrt{n}},\qquad
 r_c=\frac{\|\rho\|_2}{\|d^{(k)}\|_2},
\]
其中近零 denominator 被替换为 1；$a_c<10^{-6}$ 或
$r_c<10^{-6}$，且最多 60 coupling iterations。

\subsection{Why the spline FSI run does not converge}

对 plain spline 的 targeted echo run 在 step 62 给出决定性证据。
上一 time step 的 structural Newton 通常在 3 次 iteration 内把
residual ratio 从约 $10^{-3}$ 降到 $10^{-8}$。新 step 的第一次
structural solve 则得到：
\begin{center}
\begin{tabular}{ccc}
\toprule
Newton iteration & displacement ratio & rotation ratio\\
\midrule
1 & 55.1414 & 48.5850\\
2 & 110.237 & 103.471\\
3 & 142.074 & 104.632\\
4 & NaN & NaN\\
\bottomrule
\end{tabular}
\end{center}
因此这不是 tolerance 略微未满足，而是 full Newton step 离开
local convergence basin。

NaN 的直接机制位于 \texttt{CrBeamElement3D2N} 的 incremental
quaternion update。其 scalar increment 计算为
\[
 \Delta q_0=\sqrt{1-\frac{\|\Delta\phi\|_2^2}{4}},
\]
但 element 没有在平方根前检查 $\|\Delta\phi\|_2\le2$。当发散的
Newton correction 超过该 rotation chart 的有效域时，quaternion、
internal force 和 residual 被 NaN 污染；增加 150 iterations 或
放宽 $10^{-6}$ criterion 均不可能恢复。
受四文件修改边界限制，本次没有在 element 内插桩输出
$\Delta\phi$；因此 ``大 correction $\rightarrow$ negative square
root'' 是由 residual sequence 与该唯一显式失效路径得到的
code-level inference，而不是直接记录的 $\Delta\phi$ measurement。

触发顺序同样重要。每个新 time step 先执行
\texttt{linear\_derivative\_based} predictor，然后 solve fluid、
map reaction、solve structure；只有完成整个 coupling iterate 且
未收敛后，才对下一次 interface displacement update 应用
$\alpha=0.03$。所以 constant relaxation 无法保护新 step 的第一次
structural solve。最后一个 accepted VTK state 的 load/moment arrays
全部 finite，且最大 temporal jumps 出现在更早时刻；这排除了
``上一完整输出已经发生 load spike''，但 failed-step predictor 后的
载荷不会被 VTK flush。

结合 analytical $J^T$、basis loads 和 virtual-work tests 均通过，
该 failure 应分类为
\textbf{predictor-triggered structural nonlinear divergence at the
CrBeam incremental-rotation compatibility boundary}，而不是 mapper
algebra、factor 2 或 moment-sign error。

\subsection{Convergence design}

达到可靠收敛的顺序应为：
\begin{enumerate}
  \item 对 predicted interface displacement 使用 bounded predictor
        或 trust region，而不是无界的 full derivative extrapolation；
  \item structural solver 返回 false 或首次出现 non-finite value 时，
        立即 reject 并 rollback 当前 coupling iterate；
  \item 对 structural Newton rotation correction 做 globalization，
        保证每个 $\Delta\phi$ 留在 CrBeam quaternion chart 内；
  \item 只有完成上述 stability safeguards 后，再用 bounded Aitken
        或 IQN-ILS 恢复 coupling efficiency。
\end{enumerate}
第 2--3 项需要 solver/element support，超出当前四个 mapper 文件的
允许修改范围。因此在现有范围内应诚实报告 compatibility boundary，
不能用 parameter sweep、放宽 tolerance 或增加 iteration 上限宣称
解决。

\section{论文组织与待讨论问题}

建议 thesis 主线为：
\begin{enumerate}
  \item Ahrem generalized interpolation 与 polynomial/unisolvency；
  \item $SO(3)$ rigid-section hybrid；
  \item analytical tangent transpose 与 discrete virtual work；
  \item manufactured verification、CrBeam compatibility；
  \item Mok FSI stability evidence 与 failure classification。
\end{enumerate}

Paper 的潜在 contribution 是 finite hybrid、analytical adjoint 和
可复现 verification protocol 的组合，而不是 Rodrigues formula
本身。建议与 professor 明确讨论：
\begin{itemize}
  \item general non-commuting CrBeam motion 是否必须纳入 scope；
  \item 若必须支持，是否允许 element 提供 orientation 及其 consistent
        tangent，意味着超出当前四文件边界；
  \item FSI nonlinear structural/coupling convergence 应作为 mapper
        paper 的 acceptance gate，还是独立 numerical-solver study；
  \item hybrid 的 rigid-section assumption 是否符合目标 aerodynamic
        interface 的 physics。
\end{itemize}

\section{Reproducibility}

主要 machine-readable results 位于
\path{FSI_mok_Test_Case/TestCase_Output/MapperVerification} 和
\path{Bending_Torsion_Beam_Test_Case_2nd/TestCase_Output}。
关键 drivers 为：
\begin{itemize}
  \item \texttt{validate\_mapper\_contracts.py};
  \item \texttt{validate\_recovery\_forward.py};
  \item \texttt{validate\_recovery\_adjoint.py};
  \item \texttt{validate\_crbeam\_rotation\_compatibility.py};
  \item \texttt{validate\_bending\_torsion\_refinement.py}.
\end{itemize}

\end{document}
